\documentclass[10pt,a4paper]{amsart}
\usepackage[margin=19mm]{geometry}
\usepackage{amsmath,amssymb,mathtools}
\usepackage[scr=rsfs,cal=euler]{mathalfa}
\usepackage{xfrac}
\usepackage[unicode,hidelinks,bookmarks=false]{hyperref}
\newcommand{\ie}{i.e.\ }
\newcommand{\cf}{cf.\ }
\newcommand{\resp}{respectively\ }
\newcommand{\Subsection}{\subsection}
\providecommand{\nobreakdash}{\nobreak\hbox{-}}
\setlength{\emergencystretch}{2em}
\multlinegap0pt
\usepackage{mathtools}
\usepackage{braket}
\usepackage{csquotes}
\usepackage{bm}
\usepackage{amssymb}
\usepackage{dsfont}
\usepackage{cancel}
\usepackage{float}
\usepackage{enumerate}
\newtheorem{lemma}{Lemma}
\usepackage{cleveref}
\crefname{lemma}{Lemma}{Lemmas}
\newcommand{\HNI}{H_{\text{NI}}}
\DeclareMathOperator{\HW}{HW}
\newcommand{\cD}{\mathcal{D}}
\newcommand{\ox}{\otimes}
\title{Gibbs Sampling gives Quantum Advantage at Constant Temperatures with $O(1)$-Local Hamiltonians}
\author{Joel Rajakumar, James D. Watson}
\date{Source: arXiv:2408.01516v4 (CC BY 4.0)}
\begin{document}
\maketitle

\noindent\emph{Source and licence.} Gibbs Sampling gives Quantum Advantage at Constant Temperatures with $O(1)$-Local Hamiltonians. Version arXiv:2408.01516v4 (CC BY 4.0), distributed by arXiv under the Creative Commons Attribution 4.0 International licence, http://creativecommons.org/licenses/by/4.0/, as recorded in the arXiv licence metadata for this exact version and shown on the abstract page https://arxiv.org/abs/2408.01516v4. Full licence notice: \texttt{licenses/arxiv-2408.01516-CC-BY-4.0.txt}.\par\smallskip

\noindent\textit{Editorial scope.} Counted unit: the article's lemma Lemma:Gibbs\_to\_Circuit\_Sampling (source0.tex lines 803--809). The article's complete original proof of it is delivered with it (source0.tex lines 811--830, one \texttt{proof} environment of 20 lines). No mathematics is written, repaired or re-derived: the statement and the proof are the article's own lines, in the article's own notation, and no retained line is substituted or re-flowed.  No further source-local context is needed: every cross-reference of every retained line resolves inside the two delivered spans.  Nothing was omitted from any retained span, so the package carries no omission record.  No retained line contains a citation, so the wrapper carries no bibliography and no bibliography entry.  No retained line contains a figure environment, an image-inclusion command or drawing code, so no image is delivered and the package declares no asset dependency.  Editorial formatting only, in addition: an AMS \texttt{amsart} preamble that restates the article's own declarations and macros used by the retained lines, namely the environments \texttt{lemma} on separate counters, together with the standard packages those lines need.  The retained environments print in this document as Lemma 1 in order of appearance, whereas the article prints the same items with the numbering of its own full document.  The complete native source of the article as distributed in the arXiv e-print for this exact version is bundled unchanged as \texttt{sources/163290/source0.tex}.
	\begin{lemma}\label{Lemma:Gibbs_to_Circuit_Sampling}
		For any circuit $C$ constructed from $k$-local gates of depth $d$, there exists a $O(kd)$-local parent Hamiltonian $H_C$, such that:
  \begin{align*}
      \frac{1}{Z}e^{-\beta H_C} &= C(\cD_q^{\ox n}(\ket{0}\bra{0}))C^\dagger
  \end{align*}
  for $q = \frac{e^{-\beta}}{1+e^{-\beta}}$.
	\end{lemma}

	\begin{proof}
		Let $C$ denote an arbitrary quantum circuit.
		Define a corresponding parent Hamiltonian $H_C = C\HNI C^\dagger$ with corresponding Gibbs state:
		\begin{align}
			\frac{1}{Z}e^{-\beta H_C} &= \frac{1}{Z}e^{-\beta \sum_x \lambda_x C\ket{x}\bra{x}C^\dagger} \\
			&= \frac{1}{Z} \sum_{x\in \{0,1\}^n} e^{-\beta  \HW(x) }C\ket{x}\bra{x}C^\dagger \label{Eq:Sample_Distribution}
		\end{align} 
		Note that because unitary transformations preserve the spectrum, the partition functions of $\HNI$ and $H_C$ are the same. Let $|x|$ be hamming weight of bitstring $x \in \{0,1\}^n$.
        \noindent We can then compare this to the output state of $C$ with a single set of bit-flip noise of strength $q$ occurring on the input.
        \begin{align}
            C(\cD_q^{\ox n}(\ket{0}\bra{0}))C^\dagger &= C \left(\sum_{x \in \{0,1\}^n} q^{\HW(x)}(1-q)^{n-\HW(x)}\ket{x}\bra{x}^{\ox n} \right)C^\dagger  \\
            &= \sum_{x \in \{0,1\}^n} \bigg(\frac{q}{1-q}\bigg)^{\HW(x)}(1-q)^{n} C\ket{x}\bra{x}^{\ox n} C^\dagger \label{Eq:Noisy_Circuit}
        \end{align}
        Comparing \cref{Eq:Sample_Distribution} to \cref{Eq:Noisy_Circuit}, and noting that $Z = (1+e^{-\beta})^n$, we can see that if we set $\frac{q}{1-q} = e^{-\beta}$ and $Z = (1+e^{-\beta})^n = (1+\frac{q}{1-q})^n = 1/(1-q)^n$, then these two states are the same.
        From here it is straightforward to see that the distribution $ P(s) \coloneqq \bra{s}\frac{e^{-\beta H_C}}{Z}\ket{s}$ satisfies:
        \begin{align*}
            P(s) = P_{C, q }(s)  
        \end{align*}
        for  $q = \frac{e^{-\beta}}{1+e^{-\beta}}$.
	\end{proof}

\end{document}
